\section{Two primitive questions, some complications and a simplified matrix}\label{sec:two-questions}

\subsection{The short run and the role of money}\label{subsec:questions}

The two questions of Section~\ref{sec:io-bridge} mark key points of disjuncture in input--output modeling, and these map onto classic cleavages separating different traditions of economic thought. Such traditions can often be characterized in terms of the basic analytical ideas -- or ``rivers'' in the history of economic thought \citep[][p. 9]{Screpanti2005} -- they endorse. Two classic fields of tension in economic thought also lie at the core of paradigmatic cleavages in input--output modeling: the nature of the short run and the role of money and credit.

Labor market closures are the key entry point for different understandings of the short run: if production factors are fully utilized, increased monetary demand will shift wages and prices as capacity constraints bind. While traditionally associated with computable general equilibrium models (CGEs), this core intuition does not hold for simple models with just one production factor. Here, relative scarcity of factors does not kick in so that prices and wages stay constant and demand is simply reallocated to other sectors --- the underlying trade-off materializes in a change in the composition of output but does not affect prices. However, the key intuition can be salvaged with reference to the ``no-reallocation'' model, where prices do respond even in a one factor model as capacity constraints bind on a sectoral, instead of aggregate, level.

Conversely, if free factors of production are available, additional expenditures will lead to an expansion of employment and, correspondingly, output. This classic output adjustment closure \citep{Taylor1990} is what one gets from an application of the Leontief inverse. Both sit in one analytical apparatus \citep{Robinson.2006}, differing only in what is exogenous when monetary expenditure rises (Section~\ref{sec:kernel}). The two poles are the standard CGE closure, in which a flexible wage clears the labor market at full employment, and the employment-adjustment closure, in which a pinned money wage and constant costs set prices so that employment and output absorb demand (Table~\ref{tab:closure-conditions}). We use \emph{wage-pinned} for this class of closures, \emph{employment adjustment} for its column in the matrix, and the \emph{Leontief closure} for the fixed-coefficient corner; under demand-only shocks the column and the corner reproduce the Leontief-inverse multiplier exactly at every value of the substitution elasticities (Appendix~\ref{app:equivalences}). The former is the obvious neoclassical standard model \citep[see][]{Sen1963, TaylorLysy1979}; the latter carries the multiplier intuition of (Post-)Keynesian, radical and evolutionary approaches \citep{Pollin.2015,Mazzucato.2018}. 

Financial closures on the other hand encode prior beliefs about the role of our monetary system and the nature of money. Input-output models are typically treated as static, operating outside of time \citep[see however][]{Mercure.2018,Rocchi.2025,McGregor.1996}. A static flow model can carry domestic government borrowing, but only if some other sector supplies the relevant funds (see however \citealp{Duchin.1998}). Hence, while a traditional, simplified application of a Leontief inverse would suggest some external stimulus, the internal consistency of the model requires us to specify where that stimulus comes from. This implies that input--output models in their canonical interpretation are difficult to align with ideas on endogenous money creation as they lack any framework to conceptualize the intertemporal role of debt \citep{Lavoie.2022}. As a practical workaround the economy can be assumed to take on foreign debt, which simulates an exogenous inflow, in a less invasive way that does not threaten the internal consistency of the model.

Against this backdrop, refinancing can occur in three ways: by households (the repurposing of consumption, F1), via the state (refinancing through increased taxation, F2), or from abroad (F3), as an inflow that the static model cannot distinguish from a transfer and that carries no interest or repayment. Such an inflow provides an accounting analogue for diverse forms of external finance. Mapping these two dimensions onto specific modeling approaches yields the constellations of closure choices summarized by Table~\ref{tab:closure-conditions} below.


% =====================================================================
% Table 1 --- the two-axis core of the closure matrix, block-tagged.
% 2026-10-27: merged with the former overview table (archived); renders as Table 1.
% ---------------------------------------------------------------------
% Revised 2026-10-18 (ADR-0033): the financing conditions name the
% balancing variable. tau is the wage-income tax rate that finances the
% baseline government purchases g^G in every closure (kernel tau_rate);
% under F2 the programme g is added to that base. F is the net
% external transfer; B_gov = sum p g is the programme booked externally (F3).
% Block tags \textbf{labor market:} / \textbf{financing:} sit on their own
% line, so the two "financing:" labels align by construction.
% =====================================================================
\begin{table}[htbp]
\centering
\footnotesize
\setlength{\tabcolsep}{4pt}
\setlength{\parindent}{0pt}
\setlength{\parskip}{0pt}
\renewcommand{\arraystretch}{1.3}
\newcommand{\tblk}[1]{\textbf{#1:}}
\newcommand{\tcond}[1]{\par\hspace{1em}#1}
\newcommand{\outcome}[1]{{\centering\itshape #1\par}}
\begin{tabular}{@{}>{\raggedright\arraybackslash}p{2.2cm}
                  >{\raggedright\arraybackslash}p{6.2cm}
                  >{\raggedright\arraybackslash}p{6.2cm}@{}}
\toprule
 & \multicolumn{2}{c}{\textbf{Labor-market closure}} \\
\cmidrule(lr){2-3}
\textbf{Financing closure}
& \textbf{Standard CGE}\newline \scriptsize one flexible wage, labor mobile, full employment
& \textbf{Standard Leontief}\newline \scriptsize wage pinned, fixed coefficients, output absorbs \\
\midrule
\textbf{F1}\newline \scriptsize repurposing of consumption
&             {\outcome{pure reallocation}}
  \vspace{5pt}
\tblk{labor market}\tcond{$L_{i} = L_{i}^{\mathrm{cm}}(p, y_{i}, w)$, \quad $\textstyle\sum_{i} L_{i} = \bar L$}\par
  \vspace{6pt}
  \tblk{financing}\tcond{$\tau\, w\sum_{i} L_{i} = \textstyle\sum_{i} p_{i} g^{G}_{i}$, \quad $\tilde\beta_{i}$ tilted to program sectors, \quad $B_{\mathrm{gov}} = 0$}
& {\outcome{quantity adjustment}}
  \vspace{5pt}
\tblk{labor market}\tcond{$w = 1$, \quad $L_{i} = L_{i}^{\mathrm{cm}}(p, y_{i}, w)$, \quad $(\theta, \varepsilon, \sigma) \to 0^{+}$}\par
  \vspace{6pt}
  \tblk{financing}\tcond{$\tau\, w\sum_{i} L_{i} = \textstyle\sum_{i} p_{i} g^{G}_{i}$, \quad $\tilde\beta_{i}$ tilted to program sectors, \quad $B_{\mathrm{gov}} = 0$} \\
\midrule
\textbf{F2}\newline \scriptsize refinancing via taxes
& {\outcome{crowding out}}
  \vspace{5pt}
\tblk{labor market}\tcond{$L_{i} = L_{i}^{\mathrm{cm}}(p, y_{i}, w)$, \quad $\textstyle\sum_{i} L_{i} = \bar L$}\par
  \vspace{6pt}
  \tblk{financing}\tcond{$\tau\, w\sum_{i} L_{i} = \textstyle\sum_{i} p_{i}(g^{G}_{i} + g_{i})$, \quad $B_{\mathrm{gov}} = 0$}
& {\outcome{crowding out with a Haavelmo top-up}}
  \vspace{5pt}
\tblk{labor market}\tcond{$w = 1$, \quad $L_{i} = L_{i}^{\mathrm{cm}}(p, y_{i}, w)$, \quad $(\theta, \varepsilon, \sigma) \to 0^{+}$}\par
  \vspace{6pt}
  \tblk{financing}\tcond{$\tau\, w\sum_{i} L_{i} = \textstyle\sum_{i} p_{i}(g^{G}_{i} + g_{i})$, \quad $B_{\mathrm{gov}} = 0$} \\
\midrule
\textbf{F3}\newline \scriptsize external finance
& {\outcome{crowding out}}
  \vspace{5pt}
\tblk{labor market}\tcond{$L_{i} = L_{i}^{\mathrm{cm}}(p, y_{i}, w)$, \quad $\textstyle\sum_{i} L_{i} = \bar L$}\par
  \vspace{6pt}
  \tblk{financing}\tcond{$\tau\, w\sum_{i} L_{i} = \textstyle\sum_{i} p_{i} g^{G}_{i}$, \quad $B_{\mathrm{gov}} = \textstyle\sum_{i} p_{i} g_{i}$}
& {\outcome{output expansion and external debt}}
  \vspace{5pt}
\tblk{labor market}\tcond{$w = 1$, \quad $L_{i} = L_{i}^{\mathrm{cm}}(p, y_{i}, w)$, \quad $(\theta, \varepsilon, \sigma) \to 0^{+}$}\par
  \vspace{6pt}
  \tblk{financing}\tcond{$\tau\, w\sum_{i} L_{i} = \textstyle\sum_{i} p_{i} g^{G}_{i}$, \quad $B_{\mathrm{gov}} = \textstyle\sum_{i} p_{i} g_{i}$} \\
\bottomrule
\end{tabular}
\caption{The two-axis core of the closure matrix: the three financing closures
(rows) crossed with the two labor-market poles of the introduction (columns).
Each cell states the labor-market condition (the column) and the financing
condition (the row), and the outcome line states the qualitative response to
the demand-only impulse; the program cost $\sum_i p_i g_i$ is valued at each
cell's equilibrium prices. In the conditions, $\tau$ is the wage-income tax rate financing the
baseline purchases $g^{G}$ (a block unchanged by the program), $g$ is the
program bundle, $L_{i}^{\mathrm{cm}}$ is cost-minimizing labor demand, and
$\tilde\beta_{i}$ are expenditure shares tilted toward the program sectors
under F1.}
\label{tab:closure-conditions}
\end{table}


\subsection{Three complications}\label{subsec:complications}

The crosswise recombination of these closure options adds some complexity, but the underlying choice menu remains rough and, for the short run, simplistically dichotomous; further established variations address this latter point and thereby expand the matrix. A key mainstream route to greater variety in labor market outcomes is the leisure-income trade-off behind the canonical \textit{elastic labor supply model} \citep{deMeloTarr.1992}: Assuming that the substitution effect dominates, higher wages will attract additional labor so that capacity constraints do not bite with full force. In the \emph{elastic labor supply} closure the fixed labor endowment is replaced by a sector-specific labor supply curve, $L_i^{s} = \bar L_i\,[\,(w_i/\mathcal{P}(p))/(w_{i0}/\mathcal{P}_{0})\,]^{\eta_s}$, with one common elasticity across sectors, $\eta_s$ \citep{Roy.1951,Lagakos.2013,Galle.2023}. Two implications follow: labor supply becomes \emph{behavioral} rather than a simple quantity restriction or a price pin, and the single aggregate labor market becomes $N$ sectoral markets with $N$ wages.\footnote{The latter assumption is not strictly necessary: a single-wage form, with one aggregate supply curve $\sum_i L_i = \bar L\,[\,(w/\mathcal{P}(p))/(w_{0}/\mathcal{P}_{0})\,]^{\eta_s}$, is also possible. Under demand-only shocks, however, this single-wage form coincides with the standard CGE closure for every $\eta_s$ (Appendix~\ref{app:equivalences}), so the sectoral form is what makes the elasticity matter in our case of application, i.e.\ a shock with heterogeneous impact on sectors.}

% SOURCES (2) -- DONE 2026-10-24, verified by agent:
%   \citep{Roy.1951,Lagakos.2013,Galle.2023} on the sectoral supply curve, and
%   \citealp{Lagakos.2013,Galle.2023} at the two endpoints. All three keys are in
%   references.bib and resolve (0 undefined citations in the build).
%   Curve checked against the kernel: diagnostics.jl sectoral_supply_gap solves
%   log L_i^cm - log(labor_share_i) - eta_s,i*log(w_i/P) = 0, i.e. exactly
%   L_i = Lbar_i[(w_i/P)/(w_i0/P_0)]^{eta_s,i} with P the CPI and the anchors at 1
%   (Appendix B, eq:closure-beta). Index question settled 2026-10-24: the paper uses
%   ONE scalar eta_s everywhere (Table 4 prints eta_s in [0.5, 2], Table A4 the uniform
%   ladder); eta_{s,i} survives only in Appendix B's two provenance sentences, which
%   say the kernel CAN carry a sectoral vector (eta_s_vec) but does not here.
%   Galle entry: key Galle.2023, year 2023 -- matches the issue (ReStud 90.1 is
%   January 2023; Crossref: print 2023-01-12, online 2022-04-15), so no online-first
%   discrepancy is left in the printed entry.

While this modification brings greater responsiveness on the side of employment, other variations in the mainstream literature assume forms of rigidity that go even beyond the standard CGE labor market closure. A prominent recent example is given by \citet{BF2019}, who introduce a ``no-reallocation'' case as a baseline scenario to articulate the view that sectoral productivity shocks might impact on aggregate outcomes in a non-linear way. Relative to Table~\ref{tab:closure-conditions} this is not a column but
a more rigid, sectoral extension of the labor axis: it projects the standard CGE column's aggregate restriction on a sectoral level, i.e., $ L_i = \bar L_i$ (fixed employment levels for all sectors, so ``no-reallocation'') and correspondingly relies on sectorally flexible wages $w_i$. Formally, the transition from the standard CGE model to the no-reallocation model is governed by $\eta$, which represents the mobility of labor across sectors and admits two cases in our setup: $\eta=0$ for the no-reallocation case and $\eta=1$ for a full reallocation case that coincides with the standard CGE model.

In the context of our case of application -- a socio-ecological transformation that is tied to sector-specific shifts or injections in final demand -- the no-reallocation setup is most pessimistic for financial closures F2 and F3, as it assumes labor market rigidity across sectors, which implies that capacity constraints bite quickly and prices react accordingly. As an effect, sectors with strong demand shocks are then overwhelmed as workers from other sectors cannot ``reallocate'' to where capacity is needed. This result, however, resonates with the intuition that additional rigidities on labor markets will complicate transformation efforts that build on exactly such a reallocation of labor.\footnote{Rigidities in production, capital or financial markets are beyond the scope of this paper \citep[see, however,][]{BF2022, BernankeGertlerGilchrist.1999}.} While such an assumption deviates from the neoclassical baseline case, it strengthens the laissez-faire stance associated with its standard closure. 

Technically the \emph{Leontief closure} can be represented with the very same analytical apparatus as the standard CGE model (see Section~\ref{sec:kernel} and \ref{app:proofs}): with substitution elasticities taken to their Leontief limits \citep[see][]{Leontief.1936}, fully flexible production factors and a wage that is pinned rather than market-cleared \eqref{eq:closure-delta}, prices then follow from constant costs and final demand determines output. This is the corner most applied transformation studies work with, usually without emphasis. However, as indicated, the \emph{Leontief closure} can itself be conceived as corner solution of a broader set of related approaches that retain the basic idea that ``quantities adjust'', but do not necessarily commit to a Leontief production function. While this distinction does not matter for demand shocks (see Section~\ref{sec:application}), rigidity of production starts to matter when prices move (Section~\ref{sec:beyond-demand-only}). In their basic form these \textit{wage-pinned} models hold the money wage at its baseline value, $w=1$ \eqref{eq:closure-gamma}, which pins down prices and leaves the quantity margin free; the real wage is then an outcome rather than a restriction, since technology fixes it (Corollary~\ref{cor:realwage}). They are the (more general) closure through which the Keynesian multiplier intuition enters the model core (see also Section~\ref{sec:kernel}).

\subsection{A simplified matrix}\label{subsec:matrix}

Before we map these five closures in Table~\ref{tab:closure-matrix-3x5} below, it is worth pointing out some perhaps surprising overlaps that arise at characteristic parameter values under demand-only shocks, the case of our application, and that Figure~\ref{fig:flexibility-map} summarizes. First, as noted above, setting $\eta=1$ instead of $\eta=0$ makes the no-reallocation model coincide with the standard CGE model, since the sectoral level ceases to matter once labor reallocates fully. Second, the elastic labor supply model, which traces a curve of solutions rather than a single corner, reproduces the no-reallocation model exactly when labor market participants lack any flexibility ($\eta_s=0$) and approaches the \textit{employment-adjustment} closure at the opposite end ($\eta_s \to \infty$; see \citealp{Lagakos.2013,Galle.2023}). In that limit it matches the prices and the real wage of the employment-adjustment closure when adjusting for differences in the net external transfer.\footnote{The latter is endogenous in the elastic-supply model and zero in the employment-adjustment closure --- see Appendices~\ref{app:equivalences} and~\ref{app:sectoral} for further details.}

Hence, the elastic labor supply model can be said to sit at the center of the interparadigmatic bridge between neoclassical interpretations based on factor scarcity and more Keynesian multiplier-style interpretations based on extrapolation via the Leontief inverse \citep{Robinson.2006, Taylor1990}.

% =====================================================================
% tikz-closure-flexibility-map-v2.tex
% ---------------------------------------------------------------------
% v2 of the closure flexibility map. Changes vs. v1:
%   * ~20% narrower: every label takes a symmetric two-line break, so the
%     constant-size boxes shrink (2.7cm wide, pitch 3.2cm; v1: 3.4 / 4.0).
%   * a new double arrow between "employment adjustment" and "Leontief closure"
%     (the equivalence of Section 3: employment adjustment IS the Leontief
%     result once production is rigid), replacing the plain link there.
% Grid unchanged: row 1 = aggregate, row 2 = sectoral; columns = increasing
% flexibility. The elastic-labor-supply box still spans columns 2-3.
% Requires (already in the preamble): \usepackage{tikz} and
%   \usetikzlibrary{arrows.meta, positioning, calc}
% =====================================================================
\begin{figure}[htbp]
\centering
\begin{tikzpicture}[
    node distance = 1.4cm and 2.0cm,
    box/.style = {
      draw, rectangle, rounded corners=3pt,
      minimum width=2.7cm, minimum height=1.1cm,
      align=center, font=\sffamily\small,
      fill=gray!20, line width=0.5pt
    },
    widebox/.style = {
      box, minimum width=5.9cm          % spans two columns + one gap
    },
    rowlabel/.style = {font=\sffamily\small},
    axlabel/.style  = {font=\sffamily\small},
    link/.style     = {draw, thick},
    arr/.style      = {-{Stealth}, thick},
    flexarrow/.style = {{Stealth}-{Stealth}, thick}
  ]

  % ---- geometry -------------------------------------------------------
  \def\dx{3.2}    % horizontal pitch between column centres
  \def\rowy{1.9}  % vertical distance between row centres

  % ---- cells: row 1 = aggregate --------------------------------------
  \node[box] (r1c2) at (\dx,0)    {conventional\\CGE};
  \node[box] (r1c3) at (2*\dx,0)  {employment\\adjustment};
  \node[box] (r1c4) at (3*\dx,0)  {Leontief\\closure};

  % ---- cells: row 2 = sectoral ---------------------------------------
  \node[box]         (r2c1) at (0,-\rowy) {no\\reallocation};
  \node[widebox]     (ELS)  at (1.5*\dx,-\rowy) {elastic labor\\supply models};

  % ---- connector line through the left aggregate pair (left-right) ----
  \draw[link] (r1c2.east) -- (r1c3.west);

  % ---- double arrow: employment adjustment <-> Leontief closure --------
  \draw[flexarrow] (r1c3.east) -- (r1c4.west);

  % ---- double arrow: no reallocation -> conventional CGE
  %      up first, then right (one 90-degree corner)
  \draw[flexarrow] (r2c1.north) -- (0,0) -- (r1c2.west);

  % ---- short double arrows of the elastic-labor-supply box -----------
  \draw[flexarrow] (ELS.west)     -- (r2c1.east);        % to no reallocation
  \draw[flexarrow] (2*\dx,-1.35)  -- (r1c3.south);       % up to employment adjustment

  % ---- double arrow: single-wage elastic supply <-> conventional CGE --
  %      (the missing link: with ONE common wage the supply equation
  %      coincides with full employment for every elasticity)
  \draw[flexarrow] (\dx,-1.35) -- (\dx,-0.55);

  % ---- outer axes -----------------------------------------------------
  % y-axis with the two row labels (rotated to avoid overlap)
  \draw[arr] (-2.4,-2.7) -- (-2.4,1.0);
  \draw (-2.4,0)        -- (-2.22,0);
  \draw (-2.4,-\rowy)   -- (-2.22,-\rowy);
  \node[rowlabel, rotate=90] at (-2.75,0)      {aggregate};
  \node[rowlabel, rotate=90] at (-2.75,-\rowy) {sectoral};

  % x-axis: flexibility, increasing to the right
  \draw[arr] (-2.4,-3.2) -- (11.55,-3.2);
  \node[axlabel, below] at (4.58,-3.25) {flexibility};

\end{tikzpicture}
\caption{A map of closure choices in input--output modeling}
\label{fig:flexibility-map}
\end{figure}

 
The axiomatic variations behind closure choice overlap only partially \citep{Kapeller.2013}, and we aim to put the framework of Figure~\ref{fig:flexibility-map} to applied use by asking which of these assumptions makes what difference for the estimated economic effects of transformational policies. Since such estimates often serve as a baseline for further ecological effects and inform political decisions, the question reaches beyond the technical level. To answer it, we run the investment estimate of Section~\ref{sec:io-bridge} through different model specifications, observing how the conceptual distinctions drawn above matter in applied analysis and how the two closure dimensions (here: labor markets and financing) interact. Combining Figure~\ref{fig:flexibility-map} with the financing dimension yields the fifteen closure combinations of Table~\ref{tab:closure-matrix-3x5}, of which the two-pole core of Table~\ref{tab:closure-conditions} is the collapsed version and which shrink to nine distinct economies under demand-only shocks (Section~\ref{sec:application}).

% =====================================================================
% Table 3c --- the closure matrix, transposed
% ---------------------------------------------------------------------
% Revised 2026-10-18 (ADR-0033): the financing conditions name the balancing
% variables. tau is the wage-income tax rate that finances the baseline
% government purchases g^G in every closure (plus g under F2); B_gov (F3
% only) is the programme booked externally; F is the net external transfer
% (ADR-0019), endogenous on the mobile rows and zero on the fixed-wage rows.
% Household expenditure and the common definitions are stated once, in the
% caption.
% LAYOUT (2026-10-18 restyle, Table 6 guide-rule convention):
%   * a spanning "Labor-market closure" header over the five columns, thin
%     grey vertical guide rules between the columns (\tableline, as in
%     Table 6), and extra padding (\addlinespace) around every \midrule.
%   * every block carries a grey italic signifier on its OWN line, with its
%     conditions beneath it, ONE CONDITION PER LINE (\cnd, indented via
%     \leftskip so a wrapped condition stays indented): \lblk (labor
%     block), \fblk (financing block), \pblk (the fixed-coefficient production condition, used
%     inside the Leontief labor block).
%   * \lblk reserves a FIXED height \tbh and ends in a thin grey rule, so the
%     labor block occupies the same vertical space in every cell and the
%     rule plus the "financing" signifiers of a row sit on the SAME baseline
%     in all five columns. The Leontief cell's extra "production" block
%     lives inside that reserved height. \lblk sets \raggedright itself:
%     \parbox resets the column's alignment to justified.
%   * LINE BREAKS INSIDE CELLS: cells are p{} columns, where a literal \\
%     is the ROW separator, NOT a line break. Use \cnd for condition lines;
%     \sublabel (a \newline) is for the header and row labels only.
% \tbh is the single tunable: raise it if a condition grows a line. An
% "Overfull \vbox" warning in the log means it is too small.
% Block labels: labor market: the labor closure condition (the column);
% production: the substitution pin (only the Leontief column); financing:
% the financing condition (the row). British spelling ("labor") kept.
% =====================================================================
\begin{table}[htbp]
\centering
\footnotesize
\linespread{1.15}\selectfont% air between stacked conditions
\definecolor{tableline}{gray}{0.72}% guide rules, as in Table 6
\setlength{\tabcolsep}{3pt}
\setlength{\parindent}{0pt}
\setlength{\parskip}{0pt}
\renewcommand{\arraystretch}{1.3}
\def\tbh{60pt}% reserved height of the labor block: the tallest of the five
%   natural heights (the Leontief cell with its production condition, 58.6pt) plus
%   a point of slack. Measured, not guessed: see the header comment.
\newcommand{\tagline}[1]{\textcolor{black!60}{\textit{#1}}\par}
\newcommand{\sublabel}[1]{\newline{\scriptsize\textcolor{black!70}{#1}}}
\newcommand{\cnd}[1]{{\leftskip=0.3em\relpenalty=10000\binoppenalty=10000\relax #1\par}}
\newcommand{\pblk}[2]{\vspace{3pt}\tagline{#1}#2}
\newcommand{\lblk}[3]{\parbox[t][\tbh][t]{\linewidth}{\raggedright\tagline{#1}#2#3}\par
  \vspace{3pt}{\color{tableline}\hrule height 0.3pt}\vspace{3pt}}
\newcommand{\fblk}[2]{\vspace{3pt}\tagline{#1}#2}
\newcolumntype{I}{!{\color{tableline}\vrule width 0.3pt}}% guide rule between columns
\begin{tabular}{@{}>{\raggedright\arraybackslash}p{1.6cm}I
                  >{\raggedright\arraybackslash}p{2.51cm}I
                  >{\raggedright\arraybackslash}p{2.51cm}I
                  >{\raggedright\arraybackslash}p{2.51cm}I
                  >{\raggedright\arraybackslash}p{2.51cm}I
                  >{\raggedright\arraybackslash}p{2.51cm}@{}}
\toprule
& \multicolumn{5}{c@{}}{\textbf{Labor-market closure} \scriptsize(columns)} \\
\cmidrule(l){2-6}
\textbf{Financing closure}\sublabel{(rows)}
& \textbf{No-reallocation}\sublabel{immobile labor, sectoral wages ($\eta = 0$)}
& \textbf{Standard CGE}\sublabel{one flexible wage, mobile labor ($\eta = 1$)}
& \textbf{Elastic labor supply}\sublabel{behavioral supply on the real wage}
& \textbf{Employment adjustment}\sublabel{wage pinned, employment uncapped, $F\equiv0$}
& \textbf{Leontief closure}\sublabel{wage pinned, quantities absorb, $F\equiv0$} \\
\addlinespace[2pt]
\midrule
\addlinespace[4pt]
% ---- F1 -------------------------------------------------------------
\textbf{F1}\sublabel{repurposing of consumption}
%
& \lblk{labor market}{\cnd{$\eta = 0$:}\cnd{$L_i^{\mathrm{cm}}(p, y_i, w_i)$}\cnd{$\quad= \bar L_i$}\cnd{sectoral $w_i$}}{}
  \fblk{financing}{\cnd{$\tau \sum_i w_i L_i$}\cnd{$\quad= \sum_i p_i g^{G}_i$}\cnd{$\tilde\beta_i$ tilted to program sectors}\cnd{$B_{\mathrm{gov}} = 0$}}
%
& \lblk{labor market}{\cnd{$L_i = L_i^{\mathrm{cm}}(p, y_i, w)$}\cnd{$\sum_i L_i = \bar L$}\cnd{one wage $w$}}{}
  \fblk{financing}{\cnd{$\tau\, w \sum_i L_i$}\cnd{$\quad= \sum_i p_i g^{G}_i$}\cnd{$\tilde\beta_i$ tilted to program sectors}\cnd{$B_{\mathrm{gov}} = 0$}}
%
& \lblk{labor market}{\cnd{$L_i^{\mathrm{cm}}(p, y_i, w_i)$}\cnd{$= \bar L_i \bigl(\tfrac{w_i/\mathcal{P}(p)}{w_{i0}/\mathcal{P}_{0}}\bigr)^{\eta_s}$}}{}
  \fblk{financing}{\cnd{$\tau \sum_i w_i L_i$}\cnd{$\quad= \sum_i p_i g^{G}_i$}\cnd{$\tilde\beta_i$ tilted to program sectors}\cnd{$B_{\mathrm{gov}} = 0$}}
%
& \lblk{labor market}{\cnd{$w = 1$ (numeraire)}\cnd{$L_i = L_i^{\mathrm{cm}}(p, y_i, w)$}}{}
  \fblk{financing}{\cnd{$\tau\, w \sum_i L_i$}\cnd{$\quad= \sum_i p_i g^{G}_i$}\cnd{$\tilde\beta_i$ tilted to program sectors}\cnd{$B_{\mathrm{gov}} = 0$}}
%
& \lblk{labor market}{\cnd{$w = 1$}\cnd{$L_i = L_i^{\mathrm{cm}}(p, y_i, w)$}}{\pblk{production}{\cnd{$(\theta, \varepsilon, \sigma) \to 0^{+}$}}}
  \fblk{financing}{\cnd{$\tau\, w \sum_i L_i$}\cnd{$\quad= \sum_i p_i g^{G}_i$}\cnd{$\tilde\beta_i$ tilted to program sectors}\cnd{$B_{\mathrm{gov}} = 0$}} \\
\addlinespace[2pt]
\midrule
\addlinespace[4pt]
% ---- F2 -------------------------------------------------------------
\textbf{F2}\sublabel{refinancing via taxes}
%
& \lblk{labor market}{\cnd{$\eta = 0$:}\cnd{$L_i^{\mathrm{cm}}(p, y_i, w_i)$}\cnd{$\quad= \bar L_i$}\cnd{sectoral $w_i$}}{}
  \fblk{financing}{\cnd{$\tau \sum_i w_i L_i$}\cnd{$= \sum_i p_i(g^{G}_i + g_i)$}\cnd{$B_{\mathrm{gov}} = 0$}}
%
& \lblk{labor market}{\cnd{$L_i = L_i^{\mathrm{cm}}(p, y_i, w)$}\cnd{$\sum_i L_i = \bar L$}\cnd{one wage $w$}}{}
  \fblk{financing}{\cnd{$\tau\, w \sum_i L_i$}\cnd{$= \sum_i p_i(g^{G}_i + g_i)$}\cnd{$B_{\mathrm{gov}} = 0$}}
%
& \lblk{labor market}{\cnd{$L_i^{\mathrm{cm}}(p, y_i, w_i)$}\cnd{$= \bar L_i \bigl(\tfrac{w_i/\mathcal{P}(p)}{w_{i0}/\mathcal{P}_{0}}\bigr)^{\eta_s}$}}{}
  \fblk{financing}{\cnd{$\tau \sum_i w_i L_i$}\cnd{$= \sum_i p_i(g^{G}_i + g_i)$}\cnd{$B_{\mathrm{gov}} = 0$}}
%
& \lblk{labor market}{\cnd{$w = 1$ (numeraire)}\cnd{$L_i = L_i^{\mathrm{cm}}(p, y_i, w)$}}{}
  \fblk{financing}{\cnd{$\tau\, w \sum_i L_i$}\cnd{$= \sum_i p_i(g^{G}_i + g_i)$}\cnd{$B_{\mathrm{gov}} = 0$}}
%
& \lblk{labor market}{\cnd{$w = 1$}\cnd{$L_i = L_i^{\mathrm{cm}}(p, y_i, w)$}}{\pblk{production}{\cnd{$(\theta, \varepsilon, \sigma) \to 0^{+}$}}}
  \fblk{financing}{\cnd{$\tau\, w \sum_i L_i$}\cnd{$= \sum_i p_i(g^{G}_i + g_i)$}\cnd{$B_{\mathrm{gov}} = 0$}} \\
\addlinespace[2pt]
\midrule
\addlinespace[4pt]
% ---- F3 -------------------------------------------------------------
\textbf{F3}\sublabel{external finance}
%
& \lblk{labor market}{\cnd{$\eta = 0$:}\cnd{$L_i^{\mathrm{cm}}(p, y_i, w_i)$}\cnd{$\quad= \bar L_i$}\cnd{sectoral $w_i$}}{}
  \fblk{financing}{\cnd{$\tau \sum_i w_i L_i$}\cnd{$\quad= \sum_i p_i g^{G}_i$}\cnd{$B_{\mathrm{gov}} = \sum_i p_i g_i$}}
%
& \lblk{labor market}{\cnd{$L_i = L_i^{\mathrm{cm}}(p, y_i, w)$}\cnd{$\sum_i L_i = \bar L$}\cnd{one wage $w$}}{}
  \fblk{financing}{\cnd{$\tau\, w \sum_i L_i$}\cnd{$\quad= \sum_i p_i g^{G}_i$}\cnd{$B_{\mathrm{gov}} = \sum_i p_i g_i$}}
%
& \lblk{labor market}{\cnd{$L_i^{\mathrm{cm}}(p, y_i, w_i)$}\cnd{$= \bar L_i \bigl(\tfrac{w_i/\mathcal{P}(p)}{w_{i0}/\mathcal{P}_{0}}\bigr)^{\eta_s}$}}{}
  \fblk{financing}{\cnd{$\tau \sum_i w_i L_i$}\cnd{$\quad= \sum_i p_i g^{G}_i$}\cnd{$B_{\mathrm{gov}} = \sum_i p_i g_i$}}
%
& \lblk{labor market}{\cnd{$w = 1$ (numeraire)}\cnd{$L_i = L_i^{\mathrm{cm}}(p, y_i, w)$}}{}
  \fblk{financing}{\cnd{$\tau\, w \sum_i L_i$}\cnd{$\quad= \sum_i p_i g^{G}_i$}\cnd{$B_{\mathrm{gov}} = \sum_i p_i g_i$}}
%
& \lblk{labor market}{\cnd{$w = 1$}\cnd{$L_i = L_i^{\mathrm{cm}}(p, y_i, w)$}}{\pblk{production}{\cnd{$(\theta, \varepsilon, \sigma) \to 0^{+}$}}}
  \fblk{financing}{\cnd{$\tau\, w \sum_i L_i$}\cnd{$\quad= \sum_i p_i g^{G}_i$}\cnd{$B_{\mathrm{gov}} = \sum_i p_i g_i$}} \\
\addlinespace[2pt]
\bottomrule
\end{tabular}
\linespread{1}\selectfont% caption at the document's spacing
\caption{The expanded closure matrix: three financing closures
(rows) crossed with five labor-market closures (columns). Each cell
states the labor-market condition (the column), the fixed-coefficient production condition where
one applies (the Leontief column only), and the financing condition (the
row). The program cost $\sum_i p_i g_i$ is valued at each cell's
equilibrium prices. $\theta$, $\varepsilon$ and $\sigma$ are the substitution
elasticities whose limit $0^{+}$ gives fixed coefficients. All other notation
as in Table~\ref{tab:closure-conditions}.}
\label{tab:closure-matrix-3x5}
\end{table}


 That this space is populated rather than hypothetical is easiest to see when comparing this mapping to the macroeconometric models used in policy consulting: the EU JRC's FIDELIO \citep{Rocchi.2019, Rocchi.2025} and Cambridge Econometrics' E3ME \citep{E3ME.2022} cover both axes of the matrix, combining autonomous stimulus injections in the spirit of F3 with a wage-setting curve and involuntary unemployment that, as configured for policy work, place them in the employment-adjustment and Leontief columns, although their wage-setting modules can be reconfigured across the labor-market closures. Appendix~\ref{app:equivalences} gives the wiring; the mapping is our reading of the two model documentations, not a claim about how the models classify themselves.
